\subsection{Cobertura de identificación y dependencias excluidas}\label{sec:sd8-cobertura-identificacion}
La identificación recorre alcance y restricciones SD3/T-12, arquitectura e inventario SD4/T-11, integración y datos SD5, métodos SD6, EDT y calendarios SD7/T-14/T-15/T-18, pruebas SD9 e innovaciones SD13. El registro \path{componentes/trazabilidad_riesgos.csv} individualiza para los 34 IDs categoría causal, perspectivas T-22, fuente, paquete, horizonte y referencia al análisis. Conserva los IDs, propietarios y valores de las fichas; los 18 riesgos de innovaciones son escenarios propios vinculados a un riesgo general, no copias adicionales de la misma reserva.

La revisión del alcance detecta fallas de aceptación y capacidad en R-04/05/12; la arquitectura, obsolescencia, portabilidad y saturación en R-02/03/06/07/08; seguridad de entrega y acceso en R-09/10; ejecución física y cobertura en R-01/11/13; formación, medición y cierre en R-14/15/16. Los complementos IN1--IN5 distinguen integración, adopción, exposición documental, incentivo y atribución ambiental. La RBS clasifica estas causas en cinco grupos; cada categoría del complemento figura también en la matriz oficial de T-16. No se fuerza una fila por requisito ni se inventa una cantidad mínima de riesgos. \parencites[Formulario T-7, apartado 10; T-16; T-22]{sd8-ba}[cap.~8]{sd8-ex}

La exclusión del sistema contable no elimina el riesgo de respuesta incierta de pago o factura: R-03/R-12 requieren consulta por ID y conciliación antes de repetir un efecto. La autoridad marítima conserva la autorización de zarpe; R-02/R-03 consideran su falta de respuesta y conservan el legajo pendiente sin conceder autorización. R-06/R-11 cubren entrega de hardware, licencias y habilitación de obras/recinto del CLIENTE mediante EXT-EQUIPOS/LICENCIAS/RECINTO; R-07/R-13 cubren dependencia de proveedor y soporte. No se incluyen esos sistemas como desarrollos propios ni se supone su recepción. \parencite[cap.~11; numeral 17.1]{sd8-c07}

Cuando un evento produce efectos diferentes se conserva una respuesta por efecto: R-11 exige detener la maniobra insegura; R-10 aislar acceso a datos protegidos; R-01 suspender y reprogramar; R-02 recuperar continuidad; R-03/R-12 conciliar; R-15 mantener evidencia ambiental y no declarar certificación obtenida. R-IN4-02 protege calidad frente al incentivo, y R-IN5-01 separa atribución ambiental de continuidad de participación R-IN5-02. Un evento compartido puede activar varias respuestas, pero su suspensión y reserva se consumen una sola vez. Los 1.150 movimientos anuales de travelift son exposición operacional del caso, no 1.150 daños ni una probabilidad de falla.

\subsection{Detalle reproducible de los modos de falla}\label{sec:sd8-fmea-detalle}
Se utiliza FMEA porque permite seguir una función concreta hasta el efecto, el control y la prueba sin disponer de tasas estadísticas de falla. No se seleccionan árbol de fallas ni simulación: faltan probabilidades de eventos básicos, independencia y distribuciones verificables; no se calculan percentiles ficticios. Los cálculos físicos de sensibilidad complementan FMEA, cuyo NPR permanece ordinal. La Tabla~\ref{tab:sd8-fmea-detalle} desarrolla la función, causa y efecto de los cuatro modos ya puntuados en 8.2.
\begingroup\setlength{\tabcolsep}{3pt}
\begin{longtable}{L{.11\textwidth}L{.26\textwidth}L{.28\textwidth}L{.25\textwidth}}
\caption{Función, causa y consecuencia de los modos prioritarios}\label{tab:sd8-fmea-detalle}\\\hline
\EncabezadoTabla{ID} & \EncabezadoTabla{Función y falla} & \EncabezadoTabla{Causa y efecto} & \EncabezadoTabla{Cálculo y decisión}\\\hline\endfirsthead
\hline\EncabezadoTabla{ID} & \EncabezadoTabla{Función y falla} & \EncabezadoTabla{Causa y efecto} & \EncabezadoTabla{Cálculo y decisión}\\\hline\endhead
R-03 & Confirmar un hecho una sola vez; duplicarlo al reconectar. & Reintento, identidad o autoridad no conciliada; duplica movimiento, despacho o efecto administrativo. & $3\times5\times4=60$; probar pérdida/duplicación y bloquear efecto no conciliado.\\\hline
R-10 & Consultar sólo con finalidad y facultad; exponer salud de menor. & Permiso o vínculo incorrecto; divulgación y decisión no autorizadas. & $3\times5\times4=60$; ensayar denegación por rol/finalidad y aislar el acceso fallido.\\\hline
R-12 & Retornar sin perder hechos posteriores; volver a una imagen anterior incompleta. & Restauración sin inventario/reaplicación; descarta pagos, despachos o respuestas pendientes. & $3\times5\times3=45$; dos ensayos y conciliación de hechos posteriores antes de cortar.\\\hline
R-04 & Obtener recepción expresa antes de habilitar; cortar sin acta. & Referente ausente, agenda o expediente incompleto; activación sin recepción y atraso. & $4\times5\times2=40$; confirmar agenda/suplente y bloquear corte sin firma.\\\hline
\end{longtable}\endgroup
\FuenteTabla{Causas y efectos de las fichas T-16; escalas P/I/D declaradas por Synaptix y controles SD4/SD5/SD7.}
Detectabilidad 4 en R-03/R-10 exige prueba negativa específica, 3 en R-12 ensayo de conciliación y 2 en R-04 comprobación documental de la puerta. El orden 60/60/45/40 orienta preparación de evidencia, sin rebajar el escalamiento por impacto 5 ni demostrar eficacia ejecutada.

\subsection{Contraste de ventanas y rutas con T-18}\label{sec:sd8-contraste-ventanas}
T-18 incorpora una red global de 2.279 nodos y 5.077 enlaces, con inicio propuesto el 1 de diciembre de 2026. Su tabla de programación presenta márgenes hábiles hacia obligaciones: H3 2,33 días, H5 11, H10 14,67, corte E1 1,44 y acta/corte E2 cero. Las cuatro ramas críticas de IN1/IN3/IN4/IN5 convergen en PILOTOS-REV/PILOTOS-SUB y H12; holgura cero hacia cierre de implementación no equivale a margen hacia un hito intermedio. Estas magnitudes proceden de \path{componentes/red_global.tex} de SD7 y \path{componentes/cpm_implementacion.csv} de SD7. T-15 conserva cálculos y cargas pendientes y otro antecedente funcional: la presencia de una red en T-18 no elimina esa discrepancia ni demuestra recursos para las 34 respuestas.

La reserva marina de diez días por campaña ya imputada por T-18 no se resta otra vez a sus márgenes. Como contraste, si una suspensión total $s$ excede la duración $r$ realmente incorporada en la rama, la prolongación adicional es $a=\max(0,s-r)$ y el déficit hacia su obligación es $\max(0,a-M)$, donde $M$ es el margen de esa misma rama. Con $r=10$, $s=15$ y $M=2{,}33$, el exceso es cinco días y el déficit 2,67; con $M=11$ el déficit aritmético es cero y quedan seis días de margen. Sólo valen esas cifras si reserva, actividad y margen pertenecen a la misma cadena y no cambian ruta, ventana o recursos. No se usa margen de H5 para justificar un atraso de H3 ni para financiar el corte E2.

Si una suspensión no estaba incluida en el frente, $r=0$: la demora adicional consume su margen completo. Una demora positiva en la firma final E2, cuyo margen es cero, incumple la puerta; no la absorben diez días marinos situados antes de marcha blanca. Como la marejada suspende todos los frentes programados, las ramas concurrentes reciben el mismo bloqueo de calendario y se recalculan juntas, sin sumar cinco días por cada rama como si fueran sucesivas. Si luego comparten ejecutores, la nivelación puede aumentar la demora respecto de ese cálculo de precedencias.

Los déficits de ocho días E1/dos E2 del contraste PERT de la subred antigua son resultados de ese escenario, no fechas finales de la red global de T-18. El antecedente funcional de T-15 5.257,66/7.360,72 HH conserva su frontera; la memoria \path{componentes/estimacion_funcional.tex} de SD7 plantea una selección distinta de 12.708 HH y sensibilidad 19.640,50 HH, con otro conteo, reparto y factores. No se sustituyen silenciosamente ni se suman ambos modelos: Proyecto debe resolver la referencia conjunta y propagar la selección a cargas y red. Tampoco se presenta el incremento entre modelos como la sensibilidad de 40\,\% del antecedente de T-15.

Esta sensibilidad identifica qué hitos quedan expuestos y por qué; no constituye simulación global ni cabida acreditada de toda cancelación. Para cerrar esa cabida se requieren fechas y duraciones de cada actividad afectada, calendario protegido, frentes, recursos y reserva realmente insertada, con nueva propagación de precedencias y nivelación. Hasta conciliar T-15/T-18 se conserva la brecha y se bloquea el compromiso que dependa de ella. Los retrasos calculados son adversos condicionados; no son días esperados porque no se dispone de una probabilidad estadística del evento. \parencites[art.~17]{sd8-ba}[cap.~10, restricción 12; numeral 17.5; cap.~19]{sd8-c07}
